% Section 5.3 Testing -- Decentralized MRPII
% Uses only standard LaTeX (tabular with p-columns). Change \section/\subsection to match your report's levels.
% Test case fields follow ISO/IEC/IEEE 29119-3 (Test Documentation): ID, objective and traceability,
% preconditions, inputs, expected result, actual status.

\section{Testing}

The system was tested at three levels: unit, integration and black-box. All automated tests are written with
\textbf{pytest} and run with one command (\texttt{python -m pytest}). The test cases are documented in the style of
ISO/IEC/IEEE 29119-3: each case has an identifier, an objective traced to a design rule or decision, preconditions,
inputs, an expected result and an actual status.

Three principles shaped the test design:
\begin{itemize}
  \item \textbf{Deterministic business logic is checked with exact values.} BOM explosion, MRP, available-to-promise,
        reservations and batch arithmetic are compared with numbers worked out by hand in the design documents.
  \item \textbf{The language model is never called by a test.} Agents receive a fake model that returns chosen answers,
        broken answers or no answer, so every outcome, including the fallback, is repeatable.
  \item \textbf{Same state and same inputs give the same result.} Scenarios use a fixed seed and a simulated clock,
        and every database test starts from empty tables.
\end{itemize}

\noindent\textbf{Test environment.} Python 3 with pytest; FastAPI \texttt{TestClient} for HTTP calls; a local
PostgreSQL test database (separate from the demo database); background workers switched off and replaced by a
test ``pump'' that runs queued work step by step, so that asynchronous agent work becomes deterministic.

% ---------------------------------------------------------------------------------------------------------------------
\subsection{Unit Testing}

Unit tests check one function or module in isolation, without a database. Collaborators are replaced: the language
model by a fake client, other agents by fixed inputs. These are white-box tests: they are written with knowledge of the
code and aim at each business rule and each failure path (invalid input, refused proposal, model timeout).
140 unit tests were written; Table~\ref{tab:unit-summary} shows how they are spread over the modules.

\begin{table}[h]
\centering
\small
\renewcommand{\arraystretch}{1.2}
\begin{tabular}{|p{0.30\textwidth}|p{0.50\textwidth}|p{0.08\textwidth}|}
\hline
\textbf{Module} & \textbf{What is checked} & \textbf{Tests} \\ \hline
MRPII calculations & Multi-level MRP with capacity, ATP by date, fixed ranking of orders & 4 \\ \hline
Agent decisions (think loop) & The agent proposes an action, code checks it, one revision, then fallback & 18 \\ \hline
Core (events, rules, clock) & Event catalogue and payloads, rule bounds, batch arithmetic, configuration & 16 \\ \hline
Negotiation quotes & Quote sessions, availability steps, reading quotes back & 3 \\ \hline
API and A2A messages & Error format (RFC 9457), agent cards, refused messages & 6 \\ \hline
Purchasing Agent & Supplier selection, validation, LLM guardrails, settings and dates & 80 \\ \hline
Sales Agent & Order detail checks, Planning reply checks, customer rules & 13 \\ \hline
\multicolumn{2}{|r|}{\textbf{Total}} & \textbf{140} \\ \hline
\end{tabular}
\caption{Unit tests per module}
\label{tab:unit-summary}
\end{table}

\begin{table}[h]
\centering
\small
\renewcommand{\arraystretch}{1.2}
\begin{tabular}{|p{0.07\textwidth}|p{0.24\textwidth}|p{0.25\textwidth}|p{0.27\textwidth}|p{0.07\textwidth}|}
\hline
\textbf{ID} & \textbf{Objective (traceability)} & \textbf{Input} & \textbf{Expected result} & \textbf{Status} \\ \hline
UT-01 & Multi-level MRP with capacity is exact (Planning) & Documented phone example: BOM, stock, routings, due date
      & Gross/net requirements and planned orders equal the hand-worked numbers & Pass \\ \hline
UT-02 & ATP reacts to a supplier delay & Stock, promises, a receipt moved to a later date
      & Available-to-promise per day falls until the new receipt date & Pass \\ \hline
UT-03 & An agent's own proposal is accepted when valid (QD-7) & Fake model proposes an action with its own numbers
      & The check accepts it; the decision is recorded as chosen by the model & Pass \\ \hline
UT-04 & A wrong proposal never reaches business state (QD-11) & Fake model proposes an invalid action twice
      & One revision with the reason, then the deterministic fallback decides & Pass \\ \hline
UT-05 & A dead model does not stop the agent & Fake model times out or returns garbage
      & Fallback decides and the record states why; the loop always ends & Pass \\ \hline
UT-06 & Events are checked against the catalogue (EM-1) & Unknown, retired and wrong-publisher events
      & Each one is refused; valid payloads match their model & Pass \\ \hline
UT-07 & Rule values stay inside their bounds (RE-1) & Every business rule and its default value
      & Each default lies within its bounds; rules that depend on each other are cross-checked & Pass \\ \hline
UT-08 & Planning reads supplier lead time from quotes (DC-1) & Quotes: 60 pcs in 5 days, 500 pcs in 9 days
      & 50 pcs $\rightarrow$ 5 days; 400 pcs $\rightarrow$ 9 days; 9999 pcs $\rightarrow$ not possible & Pass \\ \hline
\end{tabular}
\caption{Selected unit test cases}
\label{tab:unit-cases}
\end{table}

% ---------------------------------------------------------------------------------------------------------------------
\subsection{Integration Testing}

Integration tests check that modules work together on the real PostgreSQL database: agents, the publish--subscribe
event channel, the private negotiation layer and the work queues. A \textbf{bottom-up} approach was used. The shared
infrastructure was tested first; then each agent on its own, with a small \emph{stub} (a fake Planning Agent that sends
requests and accepts offers); then several real agents together. The test pump plays the role of the \emph{driver}: it
runs queued work until nothing is left and fails the test if any item stays unfinished.

Besides the expected result, the agent and multi-agent tests check two global properties after each simulated day:
\begin{itemize}
  \item \textbf{Reconciliation:} each agent's records add up (stock ledger, reservations, worker hours, storage space
        and plan); the check must return no differences.
  \item \textbf{No double commitment:} no stock, space, supplier quantity or worker hour is committed twice.
\end{itemize}

\begin{table}[h]
\centering
\small
\renewcommand{\arraystretch}{1.2}
\begin{tabular}{|p{0.25\textwidth}|p{0.55\textwidth}|p{0.08\textwidth}|}
\hline
\textbf{Level} & \textbf{What is integrated} & \textbf{Tests} \\ \hline
L1 Infrastructure & Work runner (order, retries, lease take-over) and event channel (saved with the change, one copy per subscriber) & 6 \\ \hline
L2 One agent + database & Inventory, HR, negotiation layer, Purchasing Agent and Sales Agent, each against stubs & 132 \\ \hline
L3 Several agents & Purchasing with Warehouse, shortage buying, replanning, rule editor, event manager, observer & 12 \\ \hline
L4 Whole system & Full order flows started only through the simulator's public doors (see Section~5.3.3) & 8 \\ \hline
\multicolumn{2}{|r|}{\textbf{Total}} & \textbf{158} \\ \hline
\end{tabular}
\caption{Integration tests per level}
\label{tab:int-summary}
\end{table}

\begin{table}[h]
\centering
\small
\renewcommand{\arraystretch}{1.2}
\begin{tabular}{|p{0.07\textwidth}|p{0.24\textwidth}|p{0.25\textwidth}|p{0.27\textwidth}|p{0.07\textwidth}|}
\hline
\textbf{ID} & \textbf{Objective (traceability)} & \textbf{Input / steps} & \textbf{Expected result} & \textbf{Status} \\ \hline
IT-01 & A business change and its event are saved together & Publish inside a transaction, then roll back
      & Neither the change nor the event exists; on commit both exist & Pass \\ \hline
IT-02 & Failed work is retried, not lost & A handler fails repeatedly
      & Retries with back-off, then the item fails and holds its scope & Pass \\ \hline
IT-03 & Inventory keeps an exact ledger & Stub Planning reserves 3000 of 5000 pcs and commits
      & Balance, reservations and ledger rows equal the documented example & Pass \\ \hline
IT-04 & Scarce stock is never double-booked & Three requests compete for little stock
      & Total committed $\le$ stock; the others get a counter-offer or reject & Pass \\ \hline
IT-05 & HR covers a worker absence & A committed shift loses a worker
      & HR moves a multi-skilled worker, then overtime, then marks AT\_RISK & Pass \\ \hline
IT-06 & A negotiation reaches a commitment over A2A & Request $\rightarrow$ offer $\rightarrow$ accept $\rightarrow$ hold $\rightarrow$ commit
      & Session ends COMMITTED; a repeated message is a harmless duplicate & Pass \\ \hline
IT-07 & Purchasing books a receiving slot first (WH-1) & The receiving dock is full on the wanted day
      & Delivery moves to a later slot and the order still ships & Pass \\ \hline
IT-08 & A rule change reaches the agents (RE-6) & Change a rule through the rule editor
      & Agents read the new value; it survives a state restore & Pass \\ \hline
\end{tabular}
\caption{Selected integration test cases}
\label{tab:int-cases}
\end{table}

% ---------------------------------------------------------------------------------------------------------------------
\subsection{Blackbox Testing}

Black-box tests use only the system's public interface (the HTTP endpoints of Sales, Planning, the simulator and the
rule editor) and compare the outputs with the requirements, without looking at the code. The test cases were designed
with four standard techniques:
\begin{itemize}
  \item \textbf{Equivalence partitioning:} valid and invalid groups of input (a correct key and a wrong key; a complete
        order and an order with missing details).
  \item \textbf{Boundary value analysis:} rule values at, inside and outside their limits.
  \item \textbf{Decision tables:} combinations of conditions that must fit together (dependent rules).
  \item \textbf{State transition testing:} the order lifecycle, request $\rightarrow$ options $\rightarrow$ confirmation
        $\rightarrow$ commitment $\rightarrow$ production $\rightarrow$ shipment, where customer acceptance, commitment
        and production start must remain separate states.
\end{itemize}
The full scenarios (S1, S3, S8 and scenario ONE) are system-level black-box tests: they start a scenario through the
simulator and check only the visible outcome and the run report.

\begin{table}[h]
\centering
\small
\renewcommand{\arraystretch}{1.2}
\begin{tabular}{|p{0.07\textwidth}|p{0.13\textwidth}|p{0.27\textwidth}|p{0.36\textwidth}|p{0.07\textwidth}|}
\hline
\textbf{ID} & \textbf{Technique} & \textbf{Input} & \textbf{Expected result} & \textbf{Status} \\ \hline
BB-01 & Equivalence & New order request with a wrong door key / correct key
      & 403 Forbidden / request accepted & Pass \\ \hline
BB-02 & Equivalence & Order details that are wrong or missing
      & Every problem is reported; nothing is guessed & Pass \\ \hline
BB-03 & Equivalence & A2A message that is malformed, addressed to another agent, or carries a price
      & Refused with a clear error; no state changes & Pass \\ \hline
BB-04 & Boundary & Overtime rule set above its limit / within it
      & 422 ``between 0 h and 6 h'' / value applied & Pass \\ \hline
BB-05 & Decision table & Two dependent rules set so that they conflict
      & 422 ``cannot be less''; neither value is stored & Pass \\ \hline
BB-06 & State transition & Scenario S1: one normal order, customer confirms
      & States follow the lifecycle in order; the order ships complete & Pass \\ \hline
BB-07 & State transition & Cancellation after commitment
      & Exactly what was held is released; released work continues & Pass \\ \hline
BB-08 & Scenario & Scenario ONE: 150 tablets, played hour by hour
      & Planned from quotes; boards and tablets made in batches of 25; 150 shipped; records reconcile & Pass \\ \hline
BB-09 & Scenario & S3: a material shortage
      & Bought through Purchasing; the order still ships on time & Pass \\ \hline
BB-10 & Scenario & S8: four competing orders on the same day
      & No double booking; every committed order ships in full & Pass \\ \hline
BB-11 & Scenario & A supplier delay after commitment
      & Expected delivery moves; the order is marked at risk & Pass \\ \hline
BB-12 & Repeatability & Restore a saved state and replay the same day
      & The same events and the same result as the first run & Pass \\ \hline
\end{tabular}
\caption{Black-box test cases}
\label{tab:bb-cases}
\end{table}

\noindent\textbf{Test summary.} The complete suite contains 298 test functions; some run with several inputs, giving
308 test runs. In the final full run on the local test database (about 8 minutes), 307 passed, 1 was skipped and none
failed. Earlier runs had found four tests whose expectations had fallen behind the design (an old default host name, and
Planning's new make-or-buy decision record); these tests were corrected, not the system.
